% ===== PRÉAMBULE CANONIQUE — à coller VERBATIM en tête de CHAQUE .tex =====
\documentclass[11pt,a4paper]{article}
\usepackage[utf8]{inputenc}\usepackage[T1]{fontenc}\usepackage{lmodern}
\usepackage[french]{babel}
\usepackage{amsmath,amssymb,mathtools}\usepackage{icomma}
\usepackage{siunitx}\sisetup{locale=FR}  % \qty{}{}, \si{}, \ang{}
\usepackage{xcolor}\usepackage{colortbl}
\usepackage{tikz}\usepackage{pgfplots}\pgfplotsset{compat=1.18}
\usepackage[siunitx,europeanresistors]{circuitikz} % circuits (élec.)
\usepackage[most]{tcolorbox}\usepackage{enumitem}\usepackage{array,booktabs}
\usepackage[a4paper,margin=2.2cm]{geometry}
\usetikzlibrary{arrows.meta,calc,angles,quotes,positioning,patterns,%
  backgrounds,shapes.geometric,decorations.pathmorphing,decorations.markings,intersections}
\definecolor{primary}{RGB}{222,49,99}\definecolor{primarydark}{RGB}{150,22,55}
\definecolor{defcol}{RGB}{0,114,178}\definecolor{methcol}{RGB}{52,92,148}
\definecolor{excol}{RGB}{0,135,90}\definecolor{attncol}{RGB}{200,40,55}
\tcbset{boxbase/.style={enhanced,breakable,boxrule=0.9pt,arc=2.5pt,
  left=3mm,right=3mm,top=2.3mm,bottom=2.3mm,fonttitle=\bfseries,coltitle=white}}
% --- boîtes TUTORIEL (noms EXACTS, ne pas renommer) ---
\newtcolorbox{cle}[1][]{boxbase,colback=defcol!6,colframe=defcol,
  title={\textsc{À retenir}\ifx\relax#1\relax\relax\else\textmd{ : \itshape #1}\fi}}
\newtcolorbox{methode}[1][]{boxbase,colback=methcol!7,colframe=methcol,sharp corners,
  title={\textsc{Méthode}\ifx\relax#1\relax\relax\else\textmd{ --- \itshape #1}\fi}}
\newtcolorbox{exemplebox}[1][]{boxbase,colback=excol!5,colframe=excol,
  title={\textsc{Exemple}\ifx\relax#1\relax\relax\else\textmd{ : \itshape #1}\fi}}
\newtcolorbox{piege}[1][]{boxbase,colback=attncol!6,colframe=attncol,
  title={\textsc{Piège}\ifx\relax#1\relax\relax\else\textmd{ : \itshape #1}\fi}}
\newtcolorbox{remarque}[1][]{boxbase,colback=black!4,colframe=black!45,
  title={\textsc{Remarque}\ifx\relax#1\relax\relax\else\textmd{ : \itshape #1}\fi}}
% --- boîtes EXERCICE / CORRIGÉ (noms EXACTS, auto-comptées) ---
\newtcolorbox[auto counter]{exo}[1][]{boxbase,colback=primary!4,colframe=primary,
  title={\textsc{Exercice}~\thetcbcounter\ifx\relax#1\relax\relax\else\textmd{ --- \itshape #1}\fi}}
\newtcolorbox[auto counter]{corrige}[1][]{boxbase,colback=excol!4,colframe=excol,
  title={\textsc{Corrigé}~\thetcbcounter\ifx\relax#1\relax\relax\else\textmd{ --- \itshape #1}\fi}}
\newcommand{\vv}[1]{\vec{#1}}\newcommand{\ds}{\displaystyle}
\newcommand{\R}{\mathbb{R}}\newcommand{\N}{\mathbb{N}}\newcommand{\Z}{\mathbb{Z}}
\begin{document}
\sisetup{per-mode=symbol}

\begin{corrige}[trouver la position de l'image]
\[ \frac{1}{v}=\frac{1}{f}-\frac{1}{u}=\frac{1}{10}-\frac{1}{15}=\frac{3-2}{30}=\frac{1}{30}
   \ \Rightarrow\ v=+\qty{30}{\cm}. \]
$v>0$ : image \textbf{réelle}, à \qty{30}{\cm} derrière la lentille.
\end{corrige}

\begin{corrige}[grandissement et taille de l'image]
\begin{enumerate}[label=\alph*)]
  \item $\gamma=-\dfrac{v}{u}=-\dfrac{60}{20}=-3$.
  \item $i=\gamma\,o=-3\times 3=\qty{-9}{\cm}$ : image \textbf{renversée} ($i<0$) et \textbf{agrandie}
        ($|\gamma|=3$).
\end{enumerate}
\end{corrige}

\begin{corrige}[remonter à la distance focale]
\begin{enumerate}[label=\alph*)]
  \item $\dfrac{1}{f}=\dfrac{1}{u}+\dfrac{1}{v}=\dfrac{1}{30}+\dfrac{1}{60}=\dfrac{2+1}{60}=\dfrac{1}{20}
        \ \Rightarrow\ f=+\qty{20}{\cm}$.
  \item $f>0$ : lentille \textbf{convergente}.
\end{enumerate}
\end{corrige}

\begin{corrige}[lire la nature dans les signes]
Rappel : signe de $v$ $\to$ réelle/virtuelle ; signe de $\gamma$ $\to$ droite/renversée ;
$|\gamma|$ $\to$ taille.
\begin{enumerate}[label=\alph*)]
  \item \textbf{réelle} ($v>0$), \textbf{renversée} ($\gamma<0$), \textbf{réduite} ($|\gamma|<1$).
  \item \textbf{virtuelle} ($v<0$), \textbf{droite} ($\gamma>0$), \textbf{agrandie} ($|\gamma|>1$) --- loupe.
  \item \textbf{réelle}, \textbf{renversée}, de \textbf{même taille} ($|\gamma|=1$).
  \item \textbf{virtuelle}, \textbf{droite}, \textbf{réduite} --- typique d'une lentille divergente.
\end{enumerate}
\end{corrige}

\begin{corrige}[un cas complet]
\begin{enumerate}[label=\alph*)]
  \item $\dfrac{1}{v}=\dfrac{1}{20}-\dfrac{1}{60}=\dfrac{3-1}{60}=\dfrac{1}{30}\Rightarrow v=+\qty{30}{\cm}$ ;
        $\ \gamma=-\dfrac{30}{60}=-0{,}5$.
  \item $i=\gamma\,o=-0{,}5\times 2=\qty{-1}{\cm}$. Image \textbf{réelle}, \textbf{renversée},
        \textbf{réduite} de moitié (cas de l'appareil photo, $u=3f>2f$).
\end{enumerate}
\end{corrige}

\end{document}
